\subsection{Active management: value added, the information ratio and the fundamental law}

These entries measure what active management adds, relate the information ratio to the Sharpe ratio and to the attributes of a strategy, and decompose active risk. Returns are percent; ``annual'' figures from monthly data use $12\times$ the mean and $\sqrt{12}\times$ the standard deviation, which assumes serially uncorrelated monthly returns. The examples are checked in \nolinkurl{trading_notes.numerical.active_management.active_management_checks}. Forecast calibration (alpha $=$ volatility $\times$ IC $\times$ score) is in \hyperref[meth:forecast-scaling-ic]{forecast scaling}.

\begin{qfmethod}[Measuring value added: active weights and active return, tracking error, ex post and ex ante information ratios, the information ratio against the Sharpe ratio, optimal active risk and the Sharpe ratio of benchmark plus active portfolio]\label{meth:active-return-information-ratio}
\qffield{Statement} \emph{Value added} is the return relative to the benchmark. With portfolio weights $w_{Pi}$, benchmark weights $w_{Bi}$ and asset returns $r_i$, the active weights $w_{Pi}-w_{Bi}$ sum to zero and
\[ r_A=r_P-r_B=\sum_i(w_{Pi}-w_{Bi})\,r_i, \]
so a return can be computed for an active position even when its net investment is zero \citep[chap.~12 (Bailey, Richards and Tierney), p.~743, eq.~(12-9)]{managing-investment-portfolios-2007}. Active risk, or tracking error, is ``the standard deviation of active return'' \citep[chap.~3, p.~49]{alpha-grinold}. Bacon defines the ex post tracking error of arithmetic excess returns $a_i=r_i-b_i$ as $\sigma_A=\sqrt{\sum(a_i-\bar a)^2/n}$ and the information ratio as annualised excess return over annualised tracking error. He insists that the calculation be disclosed: frequency, period, arithmetic or geometric excess returns and means, $n$ or $n-1$, ex post or ex ante \citep[3rd ed., 2023, chap.~5, pp.~167--169, eqs.~(5.27)--(5.29)]{wide-bacon-attribution}.

Grinold and Kahn's information ratio is ``a ratio of (annualized) residual return to (annualized) residual risk''. Ex post it ``measures achievement''; ex ante it is ``the highest ratio of annual residual return to residual risk that the manager can obtain'' and ``connotes opportunity''. The ex post IR is approximately the $t$-statistic of alpha divided by the square root of the number of years \citep[chap.~5, pp.~110--113]{alpha-grinold}. Before fees, a top-quartile manager's IR is about 0.5 (Table~5.1, p.~114). Scaling all active positions scales alpha and residual risk together, so the IR does not depend on aggressiveness until constraints bind (pp.~115--116). It grows with the square root of the horizon: a quarterly IR is half the annual one (pp.~116--117). With value added $\alpha-\lambda_R\omega^2$, the optimal residual risk is $\omega^*=\text{IR}/(2\lambda_R)$, giving Table~5.3: 5.00 percent at IR 1.0 and $\lambda_R=0.10$, 1.67 percent at IR 0.5 and $\lambda_R=0.15$ (pp.~119--123).

\emph{Information ratio and Sharpe ratio.} Sharpe defines his ratio on a differential return $d=R_F-R_B$ against any benchmark, ex ante $\bar d/\sigma_d$ and ex post $\bar D/\sigma_D$. With the risk-free rate as benchmark it is the classic excess-return ratio; with a style benchmark it measures ``selection return''. The historical ratio equals the $t$-statistic of the mean divided by $\sqrt T$, and it is not independent of the measurement period \citep[sections ``The Ex Ante Sharpe Ratio'', ``The Ex Post Sharpe Ratio'' and ``Time Dependence'']{wide-sharpe-1994}. Interpretation: the Sharpe ratio rewards total risk and suits a portfolio held alone; the IR rewards active risk and suits a mandate measured against a benchmark.

\qffield{Notation} If the active return $A$ is uncorrelated with the benchmark's excess return, the portfolio's Sharpe ratio is $\text{SR}_P=(\mu_B+\text{IR}\cdot\sigma_A)/\sqrt{\sigma_B^2+\sigma_A^2}$. It is maximised at $\sigma_A^*=(\text{IR}/\text{SR}_B)\,\sigma_B$, where $\text{SR}_P^2=\text{SR}_B^2+\text{IR}^2$ (Interpretation, derived here; Grinold and Kahn give a relation among the benchmark's Sharpe ratio, the maximal IR and the maximal Sharpe ratio in chap.~5, Technical Appendix, item~10, p.~137). The expected active return at active risk $\sigma_A$ is $\text{IR}\cdot\sigma_A$.

\qffield{Practice} \emph{Ex post}, original synthetic year of monthly returns. The benchmark returns $2.0,-1.5,2.7,1.1,-2.4,1.2,3.5,-0.9,1.3,-2.9,2.4,0.9$; the portfolio returns $2.9,-2.1,3.1,0.9,-1.3,0.4,4.0,-0.6,0.4,-2.2,2.6,1.1$.
\begin{itemize}
\item The mean monthly active return is 0.15 percent with a sample standard deviation of 0.652 percent. Annualised, the active return is 1.80 percent, the tracking error 2.257 percent and the ex post IR 0.797.
\item Compounded, the portfolio returned 9.34 percent and the benchmark 7.40 percent: a geometric excess of 1.80 percent against an arithmetic difference of 1.94 percent.
\item With a bill rate of 0.30 percent a month, the annualised Sharpe ratios are 0.779 and 0.530.
\item The IR's $t$-statistic over one year is $0.797\times\sqrt1$, far from significant (\hyperref[meth:manager-skill-evaluation]{evaluating skill}).
\end{itemize}

\emph{Ex ante.} Take a benchmark with excess return 6 percent and standard deviation 15 percent (SR 0.40, Sharpe's illustration) and a manager with IR 0.5. The best active risk is $(0.5/0.4)\times15=18.75$ percent, giving $\text{SR}_P=\sqrt{0.16+0.25}=0.640$. A 4 percent tracking-error limit gives 0.515. Doubling active risk past the optimum, to 37.5 percent, lowers the ratio to 0.613. Grinold and Kahn's own objective, with residual-risk aversion $\lambda_R=0.10$, chooses only 2.5 percent at IR 0.5: their $\lambda_R$ is a mandate's aversion to active risk, not total-portfolio mean--variance optimisation.

\qffield{Conditions} Annualising by 12 and $\sqrt{12}$ assumes independent monthly active returns; serial correlation or smoothing makes the annualised IR wrong (\hyperref[meth:skill-luck-significance-tests]{skill against luck}). Grinold and Kahn's residual return is relative to beta times the benchmark; with portfolio beta 1, residual and active risk coincide (p.~111). An ex ante IR is a forecast of skill and is unobservable; an ex post IR is a noisy estimate.

\qffield{Pitfalls} Comparing IRs computed with different conventions. Reading a high IR as skill over short samples. Taking active risk beyond the optimum because the IR does not change with scale: real constraints, costs and the investor's total risk then dominate. Using the Sharpe ratio to compare managers who will be combined with a benchmark core, where the IR is the relevant ratio.

\qffield{See also} \hyperref[res:fundamental-law-active-management]{fundamental law}; \hyperref[meth:active-risk-decomposition]{active risk decomposition}; \hyperref[meth:risk-adjusted-performance-measures]{risk-adjusted performance measures}.

\end{qfmethod}

\begin{qfresult}[The fundamental law of active management: $\text{IR}\approx\text{TC}\cdot\text{IC}\cdot\sqrt{\text{BR}}$; breadth, skill and transfer; choosing active risk, comparing timing with selection, additivity across strategies and managers; independence and other limits]\label{res:fundamental-law-active-management}
\qffield{Statement} The fundamental law ``gives an approximation to the information ratio'' from two attributes of a strategy: \emph{breadth} BR, ``the number of independent forecasts of exceptional return we make per year'', and skill, the \emph{information coefficient} IC, ``the correlation of each forecast with the actual outcomes'', assumed equal for all forecasts:
\[ \text{IR}\approx\text{IC}\cdot\sqrt{\text{BR}} \]
\citep[chap.~6, pp.~147--148, eq.~(6.1)]{alpha-grinold}. The approximation ignores the risk reduction that knowledge brings, which is very small for IC below 0.1 (p.~148). With $\omega^*=\text{IR}/(2\lambda_R)$, the optimal residual risk and value added become $\text{IC}\sqrt{\text{BR}}/(2\lambda_R)$ and $\text{IC}^2\,\text{BR}/(4\lambda_R)$: risk rises with skill and with the square root of breadth (pp.~148--149). Clarke, de Silva and Thorley (2002) added the \emph{transfer coefficient} TC to account for imperfect implementation, $\text{IR}=\text{IC}\cdot\sqrt{\text{BR}}\cdot\text{TC}$. TC is the correlation between the optimal unconstrained ``paper'' portfolio Q and the portfolio P actually held, and $\text{IR}_P=\text{IR}_Q\cdot\rho_{PQ}$ \citep[chap.~3, pp.~23 and~27--28]{wide-grinold-kahn-advances}.

\qffield{Conditions} The law rests on three assumptions: the manager measures his skill accurately and uses information optimally, the most crucial; the sources of information are independent, so the manager ``doesn't bet twice on a repackaged form of the same information''; and the IC is the same for each forecast \citep[chap.~6, p.~162]{alpha-grinold}. Dependence lowers combined skill: two sources with IC $0.0816$ each and correlation $\gamma=0.5$ give a combined IC of 0.0942, not $\sqrt2\times0.0816$ (p.~159). Reassessing an industry view once a year but rebalancing monthly is ``the same bet 12 times'', one unit of breadth (p.~158). The law is not the law of large numbers: it holds at any breadth if skill is maintained (p.~160). It is ``a guideline, not an operational tool'', and breadth is hard to estimate because forecasts must be independent (pp.~149 and~160). Long-only portfolios with further constraints ``can experience transfer coefficients well below 0.5'', while long-short portfolios of cheap-to-trade assets can exceed 0.9 \citep[p.~29]{wide-grinold-kahn-advances}. Transaction costs also cut the realised IR: in Grinold and Kahn's binary example, 0.80 percent round-trip costs on 100 percent annual turnover lower it from 1.16 to 0.70 \citep[n.~4, p.~153]{alpha-grinold}.

\qffield{Practice} \emph{Applications} (figures reproduced from the sources).
\begin{itemize}
\item \emph{Timing against selection.} For IR 0.5, a market timer with four independent quarterly calls needs IC 0.25; a stock picker making 400 bets a year needs 0.025 (chap.~6, p.~149). In the sequel, a stock picker with IC 0.05, 2{,}000 bets and TC 0.35 reaches $0.05\times\sqrt{2000}\times0.35=0.78$. A long-only market timer with IC 0.1, four bets and TC 0.6 reaches 0.12. A tactical allocator with IC 0.1, 12 bets and TC 0.5 reaches 0.17 \citep[p.~29]{wide-grinold-kahn-advances}.
\item \emph{Adding a research area.} Squared IRs add for independent sources: 200 stocks twice a year at IC 0.04 give $0.04^2\times400=0.64$; adding 100 stocks at IC 0.03 adds $0.03^2\times200=0.18$, raising value added 28 percent and the IR from 0.80 to 0.91 \citep[chap.~6, p.~154]{alpha-grinold}.
\item \emph{Combining managers.} Three independent managers with IRs 0.75, 0.50 and 0.30, optimally allocated, give $\sqrt{0.75^2+0.5^2+0.3^2}=0.95$ (p.~157).
\item \emph{Choosing active risk and forecasting value added.} With IR known, expected active return is $\text{TC}\cdot\text{IC}\cdot\sqrt{\text{BR}}\cdot\sigma_A$: 3.13 percent at $\sigma_A=4$ percent for the stock picker above.
\item \emph{A check.} In the binary example (residual returns made of 300 equally likely $\pm1$ percent components, the manager knowing one; 200 stocks; two forecasts a year), $\text{IC}=1/\sqrt{300}=0.0577$ and the law predicts $0.0577\times\sqrt{400}=1.155$. The exact portfolio has active standard deviation 1.2227 percent per half-year, IR 0.818 per half-year and 1.157 a year (pp.~151--153).
\end{itemize}
\emph{Transfer coefficient, original.} Five stocks have residual risks 20, 25, 30, 20 and 25 percent and benchmark-neutral alphas 1.0, 0.5, $-0.5$, 1.0 and $-2.0$ percent. The unconstrained optimum ($\lambda=0.005$) wants active positions of 25, 8, $-5.6$, 25 and $-32$ points, with $\text{IR}_Q=0.110$. Limits of $-5$ to $+10$ points bind on four stocks. The constrained portfolio has $\text{TC}=0.835$ and $\text{IR}_P=0.835\times0.110=0.092$ exactly, as the identity requires.

\qffield{Pitfalls} Counting holdings or rebalancing dates as breadth. Assuming IC is stable across stocks and periods. Expecting the unconstrained IR from a long-only, capped portfolio. Using the law to promise performance: IC and BR are forecasts. The law encourages diversified, independent sources (p.~161); a sponsor diversifying across managers gains breadth only if their bets are independent.

\qffield{See also} \hyperref[meth:active-return-information-ratio]{value added and IR}; \hyperref[meth:forecast-scaling-ic]{forecast scaling}; \hyperref[meth:alpha-signal-combination]{signal combination}.

\end{qfresult}

\begin{qfmethod}[Sources of active risk: active factor exposures against active specific risk, active beta against residual risk, marginal contributions, and the interpretation of tracking risk and the information ratio]\label{meth:active-risk-decomposition}
\qffield{Statement} Active risk comes from active exposures to common factors and from active positions in individual stocks. With factor exposures $X$, factor covariance $F$ and specific variances $\Delta$, the active holdings $h_A=h_P-h_B$ have active exposures $x_A=X^\top h_A$, and because factor and specific returns are uncorrelated,
\[ \psi^2=x_A^\top F x_A+h_A^\top\Delta h_A, \]
separating ``both total and active risk into common-factor and specific components'' \citep[chap.~3, Technical Appendix, pp.~76--77]{alpha-grinold}. Relative to the benchmark, the same variance splits into active systematic and residual parts, $\psi_P^2=\beta_{PA}^2\sigma_B^2+\omega_P^2$, with active beta $\beta_{PA}=\beta_P-1$ (chap.~4, p.~90). Residual risk is ``the risk of the return orthogonal to the systematic return'' (chap.~3, p.~50). Marginal contributions to active risk, $\partial\psi/\partial h_A=Vh_A/\psi$, weighted by the active holdings, sum to $\psi$, and they identify the positions that diversify or concentrate active risk (p.~67).

\qffield{Practice} Example, original. Five stocks with exposures to three factors (market, size, value), factor volatilities 16, 8 and 6 percent (correlations 0.3, 0.1 and 0.2) and specific volatilities 30, 25, 35, 28 and 32 percent. The active weights are $+8$, $-4$, $+5$, $-6$ and $-3$ points against an equal-weight benchmark.
\begin{itemize}
\item The active exposures are 0.039 to the market, 0.131 to size and $-0.137$ to value.
\item Active variance is 15.675 percent squared, so the tracking error is 3.96 percent. Factor variance is 2.109 (13.5 percent of the total) and specific variance 13.567.
\item Active beta is 0.045, contributing only 0.93 percent squared; residual variance is 14.75.
\item The holdings-weighted marginal contributions add to 3.96 percent.
\end{itemize}
Interpretation: this is a stock-selection book. Its active risk is mostly specific, so its IR depends on stock-level skill and breadth, and a sponsor wanting factor exposure should get it more cheaply elsewhere. A book whose active variance were mostly size and value would be a factor bet, judged against a factor benchmark (\hyperref[meth:benchmark-selection-quality]{benchmark selection}).

\qffield{Conditions} The decomposition is only as good as the factor model and covariance estimates (\hyperref[mod:factor-risk-model]{factor risk model}). Ex ante tracking error from a model differs from ex post tracking error measured on returns, and the two should be labelled \citep[p.~167]{wide-bacon-attribution}.

\qffield{Pitfalls} Reading low tracking error as good management; Bacon notes that it ``might imply consistent underperformance'' (p.~167). Allocating covariance between correlated factors as if it belonged to one of them. Ignoring that an active beta away from zero makes the manager a benchmark timer.

\qffield{See also} \hyperref[mod:factor-risk-model]{factor risk model}; \hyperref[def:factor-model-types]{factor model types}; \hyperref[meth:active-return-information-ratio]{value added and IR}.

\end{qfmethod}
